
\documentclass[11pt]{article}
\usepackage[utf8]{inputenc}
\usepackage{amsmath, amssymb}
\usepackage{braket}
\usepackage{quantikz}
\usepackage{geometry}
\usepackage{graphicx}
\geometry{a4paper, margin=1in}

\title{Quantum Phase Estimation (QPE)}
\author{Morten HJ}
\date{\today}

\begin{document}

\maketitle
\tableofcontents
\newpage

% Introduction
\section{Introduction}
Quantum Phase Estimation (QPE) is a fundamental quantum algorithm that estimates the eigenphase $\phi$ of a unitary operator $U$. Given an eigenstate $\ket{\psi}$ satisfying:
\begin{equation}
    U\ket{\psi} = e^{2\pi i \phi} \ket{\psi},
\end{equation}
QPE provides an estimate of $\phi$ with high probability. It is a crucial subroutine for algorithms like Shor's factoring algorithm and quantum simulations.

% Mathematical Background
\section{Mathematical Background}
The goal of QPE is to estimate the phase $\phi \in [0,1)$ where:
\begin{equation}
    U\ket{\psi} = e^{2\pi i \phi} \ket{\psi}.
\end{equation}

\subsection{Quantum Fourier Transform (QFT)}
The Quantum Fourier Transform (QFT) on $n$ qubits is defined as:
\begin{equation}
    \ket{x} \rightarrow \frac{1}{\sqrt{2^n}}\sum_{y=0}^{2^n-1} e^{2\pi i xy / 2^n}\ket{y}.
\end{equation}
QFT is essential for extracting phase information in QPE.

\newpage

% Quantum Circuit and Working
\section{Quantum Circuit and Working}
The quantum circuit for QPE consists of two quantum registers:
\begin{itemize}
    \item The first register with $n$ qubits is initialized to $\ket{0}^{\otimes n}$.
    \item The second register is initialized to the eigenstate $\ket{\psi}$ of $U$.
\end{itemize}

\begin{figure}[h!]
    \centering
    \begin{quantikz}
        \lstick{$\ket{0}^{\otimes n}$} & \gate{H} & \ctrl{1} & \gate{QFT^\dagger} & \meter{} \\
        \lstick{$\ket{\psi}$} & \qw & \gate{U^{2^j}} & \qw & \qw
    \end{quantikz}
    \caption{Quantum Phase Estimation Circuit}
\end{figure}

% Derivation of the Algorithm
\section{Derivation of the Algorithm}
\subsection{Step 1: Initialization}
The system starts in the state:
\begin{equation}
    \ket{0}^{\otimes n} \otimes \ket{\psi}.
\end{equation}
Applying Hadamard gates to the first register creates a superposition:
\begin{equation}
    \frac{1}{2^{n/2}}\sum_{k=0}^{2^n -1} \ket{k}\otimes \ket{\psi}.
\end{equation}

\subsection{Step 2: Controlled Unitary Operations}
Controlled-$U^{2^j}$ gates apply the operation $U$ with exponential powers:
\begin{equation}
    \frac{1}{2^{n/2}}\sum_{k=0}^{2^n -1} \ket{k} \otimes U^k\ket{\psi}.
\end{equation}
For eigenstate $\ket{\psi}$, this becomes:
\begin{equation}
    \frac{1}{2^{n/2}}\sum_{k=0}^{2^n -1} e^{2\pi i \phi k}\ket{k} \otimes \ket{\psi}.
\end{equation}

\subsection{Step 3: Quantum Fourier Transform (QFT)}
Applying QFT to the first register transforms the state to:
\begin{equation}
    \sum_{k=0}^{2^n -1} c_k \ket{k},
\end{equation}
where coefficients $c_k$ are peaked near $k \approx 2^n \phi$.

\subsection{Step 4: Measurement}
Measuring the first register gives an $n$-bit estimate of $\phi$ with high probability.

\newpage

% Applications
\section{Applications of QPE}
QPE is a foundational algorithm with several key applications:
\begin{itemize}
    \item \textbf{Factoring and Cryptography:} Integral to Shor's algorithm for integer factoring.
    \item \textbf{Quantum Simulations:} Estimating eigenvalues of Hamiltonians in quantum chemistry.
    \item \textbf{Amplitude Estimation:} Enhances algorithms like Grover's search.
\end{itemize}

% Challenges and Practical Considerations
\section{Challenges and Practical Considerations}
\begin{itemize}
    \item \textbf{Qubit Requirements:} Requires a large number of qubits for high precision.
    \item \textbf{Gate Depth:} Controlled-$U^{2^j}$ operations increase gate complexity.
    \item \textbf{Noise Sensitivity:} Errors in QFT or controlled gates impact accuracy.
\end{itemize}

% Conclusion
\section{Conclusion}
Quantum Phase Estimation is a powerful quantum algorithm for extracting phase information of unitary operators. It forms the foundation for many advanced quantum algorithms and demonstrates the power of quantum Fourier transforms in computational tasks.

% References
\section{References}
\begin{thebibliography}{9}
\bibitem{nielsen} M. A. Nielsen and I. L. Chuang, \textit{Quantum Computation and Quantum Information}, Cambridge University Press, 2000.
\bibitem{preskill} J. Preskill, \textit{Lecture Notes on Quantum Computation}, Caltech.
\end{thebibliography}

\end{document}
